\subsection*{Logbook Entry 31: Focused candidate evidence and regulatory-overlap follow up}

\textbf{What was being measured.}
A focused follow up was performed on the six candidate genes detected in the full-universe D3 versus not-D3 methylation scan, with the aim of comparing their probe-level evidence, regulatory annotations, and relationship to the original top-200 directional export. The broader question was whether the SIM2-centered 21q signature sits inside a more coherent oncogenic developmental program.

\textbf{Methods.}
Probe-level rows were extracted from the full-universe D3 versus not-D3 table for PAX6, GLI2, BMP4, PTCH1, SHH, and OLIG2, after manifest-based annotation merge. For each candidate, the analysis summarized probe count, best absolute median difference, minimum adjusted \(q\)-value, chromosome, genomic context, and available regulatory annotations including regulatory feature group, regulatory feature name, PHANTOM4 enhancer labels, DNase labels, open chromatin labels, and TFBS labels. The same candidate set was then intersected with the original top-200 directional methylation export to determine which signals rose to the shortlist level.

\textbf{Results.}
The full-universe candidate set contained probe-level support for PAX6, GLI2, BMP4, PTCH1, SHH, and OLIG2. Shared regulatory feature groups across these genes were: Unclassified\_Cell\_type\_specific. Intersection with the original top-200 directional export showed top-200 presence for: PAX6. This establishes which components of the broader developmental bridge were strong enough to enter the original shortlist and which remained weaker background candidates visible only in the broader universe scan.

\textbf{Interpretation.}
This step helps determine whether the SIM2-centered signal can reasonably be framed as part of a broader oncogenic developmental program. The logic is not merely nominal overlap of gene names, but convergence across methylation effect size, regulatory-context annotation, and shortlist prominence. Shared regulatory annotations such as enhancer, TFBS, DNase, or open-chromatin labels would strengthen the case for coordinated regulatory structure, whereas scattered weak hits without shared regulatory context would argue for a looser developmental backdrop rather than a coherent program.

\textbf{Tables written.}
\texttt{results/gse240704/focused\_candidate\_evidence/focused\_candidate\_full\_universe\_hits.csv}\\
\texttt{results/gse240704/focused\_candidate\_evidence/focused\_candidate\_summary.csv}\\
\texttt{results/gse240704/focused\_candidate\_evidence/focused\_candidate\_regulatory\_overlap.csv}\\
\texttt{results/gse240704/focused\_candidate\_evidence/focused\_candidate\_top200\_intersection.csv}\\
\texttt{results/gse240704/focused\_candidate\_evidence/tables\_tex/table\_focused\_candidate\_evidence.tex}\\
\texttt{results/gse240704/focused\_candidate\_evidence/tables\_tex/table\_focused\_candidate\_top200\_intersection.tex}\\
\texttt{results/gse240704/focused\_candidate\_evidence/focused\_candidate\_summary.json}

\textbf{Figures.}
\texttt{results/gse240704/focused\_candidate\_evidence/plots/focused\_candidate\_best\_effect.png}

\textbf{Next steps.}
The next step is to use the focused candidate evidence to decide whether the manuscript should describe the broader background as a coordinated developmental program or as a weaker distributed backdrop around the dominant local SIM2 signal. If regulatory-context sharing is sparse, interpretation should remain cautious. If shared enhancer, TFBS, DNase, or shortlist enrichment is clear across several genes, the developmental-program language becomes more defensible.
